\par\medskip\noindent\textbf{Cómo leer la salida.}
{\small\setlength{\tabcolsep}{4pt}\renewcommand{\arraystretch}{1.15}
\begin{xltabular}{\linewidth}{@{}>{\raggedright\arraybackslash}p{0.24\linewidth}>{\raggedleft\arraybackslash}p{0.15\linewidth}>{\raggedright\arraybackslash}X@{}}
\toprule Indicador & Valor & Interpretación \\ \midrule \endfirsthead
\toprule Indicador & Valor & Interpretación \\ \midrule \endhead
\bottomrule \endlastfoot
$P(X\le4)$ & 0,8922 & Con 250 días y VaR correcto, hasta 4 excepciones ocurren el 89\,\% de las veces: zona verde. \\
$P(X\le9)$ & 0,9997 & 10 o más excepciones son casi imposibles si el modelo es correcto: zona roja. \\
Columnas por año & p. ej. 2008: 21 roja & Número de excepciones y zona de Basilea (verde 0--4, amarilla 5--9, roja 10 o más). \\
2007 y 2008 & roja en HS y EWMA & Los años de crisis concentran los problemas. \\
2014 EWMA & 10 roja & La EWMA con normal falla incluso en un año sin crisis por subestimar las colas. \\
Años en verde / amarilla / roja & HS 12/4/2; EWMA 8/8/2; GARCH-$t$ 12/6/0 & El GARCH-$t$ nunca cae en zona roja. \\
Multiplicador medio & 3,23 / 3,37 / 3,20 & Factor por el que el regulador multiplica el VaR para el capital (3 en verde, hasta 4 en roja); el GARCH-$t$ exige menos capital. \\
\end{xltabular}}
